\originalchapter{8.333期末复习包：七个历史年份的量子与相变题}
\osource{Fall 2013页面的2007 Final Review及官方Solutions，内含Fall1998、1999、2000、2003、2004、2005、2006各三题，共21个原编号。没有2001、2002卷，不补造；这些是包内标明年份的历史考试，不是2013年实际期末卷。}
\originalsection{1998年：激子与离子凝聚}
\begin{exercise}[1998题1(a)--(e)]\label{f:333-1998-1}
半导体光激发$N_e$电子、$N_h$空穴，质量$m_e,m_h$，可结合成质量$M=m_e+m_h$、降低能$\varepsilon$的激子，均按经典非相互作用气、忽略内部简并。(a)自由电子空穴自由能；(b)$N_p$激子自由能；(c)三化学势；(d)平衡条件；(e)高温、总自由激发密度$n\simeq n_e+n_h$时的激子密度，取电中性。
\end{exercise}
\begin{solution}
$\lambda_j=h/\sqrt{2\pi m_j\kb T}$，$F_e+F_h=\kb T\sum_{j=e,h}N_j[\log(n_j\lambda_j^3)-1]$，$F_p=N_p\{\kb T[\log(n_p\lambda_p^3)-1]-\varepsilon\}$。所以$\mu_e=\kb T\log(n_e\lambda_e^3)$、$\mu_h$类同、$\mu_p=-\varepsilon+\kb T\log(n_p\lambda_p^3)$。

电子加空穴变一激子，平衡$\mu_e+\mu_h=\mu_p$，即$n_p=n_en_h(\lambda_e^3\lambda_h^3/\lambda_p^3)\e^{\beta\varepsilon}=n_en_h\lambda_\mu^3\e^{\beta\varepsilon}$，$\mu=m_em_h/(m_e+m_h)$为约化质量。高温解离多，$n_e=n_h\simeq n/2$，$n_p\simeq n^2\lambda_\mu^3\e^{\beta\varepsilon}/4$。更低温需用激发对守恒$n_{pair}=n_e+n_p=n_h+n_p$解二次方程，不能保持$n_e=n/2$近似到全结合区。
\end{solution}
1998题2(a)--(e)Manning题逐项在既有习题33.3。此历史题用固定$\log(r/L)$而PS5用$\log(r/R)$：只把单粒子$z$的前因子由$R^{2\xi}$换成$L^{2\xi}$，径向概率、均值、凝聚反馈不变，物理壁压按固定能零接触式计算。1998题3(a)--(g)逐项对应本卷习题\ref{p:333-6-11}；热容全温峰与二维极限的修正同样适用，官方三维临界峰公式不能外推成所有维数的全温最大值。
\originalsection{1999年：Fermi与Bose磁性}
\begin{exercise}[1999题1(a)--(e)]\label{f:333-1999-1}
密度$n=N/V$的spin1/2自由电子。(a)基态、$k_F,E_0/V$；加入$U=\alpha N_+N_-/V$、$\alpha=4\pi\hbar^2a/m$模拟偏好平行自旋。(b)两Fermi波数；(c)$n_\pm=n/2\pm\delta$的小偏极动能；(d)相互作用与Stoner$\alpha_c$；(e)自发磁化随$\alpha$及图。
\end{exercise}
\begin{solution}
基态每自旋占满动量球，$k_F=(3\pi^2n)^{1/3}$、$E_F=\hbar^2k_F^2/(2m)$，$E_0/V=3nE_F/5$。偏极$k_{F\pm}=(6\pi^2n_\pm)^{1/3}$，动能密度$\tfrac35[\hbar^2(6\pi^2)^{2/3}/(2m)](n_+^{5/3}+n_-^{5/3})$，至$\delta^2$增加$4E_F\delta^2/(3n)$。相互作用$\alpha(n^2/4-\delta^2)$，故$\alpha_c=4E_F/(3n)$。

用$m_s=2\delta/n$，能至四阶$E/V=\mathrm{const}+(nE_F/3)(1-\alpha/\alpha_c)m_s^2+(nE_F/81)m_s^4+\cdots$，近临界$m_s\simeq\sqrt{(27/2)(\alpha/\alpha_c-1)}$。大耦合需最小化完整$5/3$次式，磁化最后饱和$|m_s|=1$；饱和边界$\alpha/\alpha_c=3\,2^{2/3}/4\simeq1.19055$。下图用完整方程，未将小磁化展开外推到饱和值。
\begin{center}\begin{tikzpicture}\begin{axis}[width=10cm,height=4.5cm,xlabel={$\alpha/\alpha_c$},ylabel={$|m_s|$},xmin=.8,xmax=1.4,ymin=0,ymax=1.05]\addplot[thick,structurecolor] coordinates {(0.8,0) (0.803351955,0) (0.806703911,0) (0.810055866,0) (0.813407821,0) (0.816759777,0) (0.820111732,0) (0.823463687,0) (0.826815642,0) (0.830167598,0) (0.833519553,0) (0.836871508,0) (0.840223464,0) (0.843575419,0) (0.846927374,0) (0.85027933,0) (0.853631285,0) (0.85698324,0) (0.860335196,0) (0.863687151,0) (0.867039106,0) (0.870391061,0) (0.873743017,0) (0.877094972,0) (0.880446927,0) (0.883798883,0) (0.887150838,0) (0.890502793,0) (0.893854749,0) (0.897206704,0) (0.900558659,0) (0.903910615,0) (0.90726257,0) (0.910614525,0) (0.91396648,0) (0.917318436,0) (0.920670391,0) (0.924022346,0) (0.927374302,0) (0.930726257,0) (0.934078212,0) (0.937430168,0) (0.940782123,0) (0.944134078,0) (0.947486034,0) (0.950837989,0) (0.954189944,0) (0.957541899,0) (0.960893855,0) (0.96424581,0) (0.967597765,0) (0.970949721,0) (0.974301676,0) (0.977653631,0) (0.981005587,0) (0.984357542,0) (0.987709497,0) (0.991061453,0) (0.994413408,0) (0.997765363,0) (1.00111732,0.12245661) (1.00446927,0.242773006) (1.00782123,0.318360244) (1.01117318,0.377207494) (1.01452514,0.42635704) (1.01787709,0.468914866) (1.02122905,0.50658661) (1.02458101,0.540432723) (1.02793296,0.571168223) (1.03128492,0.599303719) (1.03463687,0.625219856) (1.03798883,0.64921001) (1.04134078,0.671506405) (1.04469274,0.692296914) (1.04804469,0.711736305) (1.05139665,0.729954045) (1.0547486,0.74705987) (1.05810056,0.763147843) (1.06145251,0.7782994) (1.06480447,0.792585657) (1.06815642,0.806069196) (1.07150838,0.818805468) (1.07486034,0.830843905) (1.07821229,0.842228815) (1.08156425,0.85300011) (1.0849162,0.863193903) (1.08826816,0.872843001) (1.09162011,0.881977314) (1.09497207,0.890624206) (1.09832402,0.89880878) (1.10167598,0.906554133) (1.10502793,0.913881562) (1.10837989,0.92081075) (1.11173184,0.927359921) (1.1150838,0.933545975) (1.11843575,0.939384606) (1.12178771,0.944890405) (1.12513966,0.950076942) (1.12849162,0.954956849) (1.13184358,0.959541879) (1.13519553,0.963842963) (1.13854749,0.967870256) (1.14189944,0.971633169) (1.1452514,0.975140398) (1.14860335,0.97839994) (1.15195531,0.981419091) (1.15530726,0.984204443) (1.15865922,0.986761844) (1.16201117,0.989096351) (1.16536313,0.991212125) (1.16871508,0.993112282) (1.17206704,0.994798638) (1.17541899,0.99627127) (1.17877095,0.997527735) (1.18212291,0.99856146) (1.18547486,0.999357859) (1.18882682,0.999880604) (1.19217877,1) (1.19553073,1) (1.19888268,1) (1.20223464,1) (1.20558659,1) (1.20893855,1) (1.2122905,1) (1.21564246,1) (1.21899441,1) (1.22234637,1) (1.22569832,1) (1.22905028,1) (1.23240223,1) (1.23575419,1) (1.23910615,1) (1.2424581,1) (1.24581006,1) (1.24916201,1) (1.25251397,1) (1.25586592,1) (1.25921788,1) (1.26256983,1) (1.26592179,1) (1.26927374,1) (1.2726257,1) (1.27597765,1) (1.27932961,1) (1.28268156,1) (1.28603352,1) (1.28938547,1) (1.29273743,1) (1.29608939,1) (1.29944134,1) (1.3027933,1) (1.30614525,1) (1.30949721,1) (1.31284916,1) (1.31620112,1) (1.31955307,1) (1.32290503,1) (1.32625698,1) (1.32960894,1) (1.33296089,1) (1.33631285,1) (1.3396648,1) (1.34301676,1) (1.34636872,1) (1.34972067,1) (1.35307263,1) (1.35642458,1) (1.35977654,1) (1.36312849,1) (1.36648045,1) (1.3698324,1) (1.37318436,1) (1.37653631,1) (1.37988827,1) (1.38324022,1) (1.38659218,1) (1.38994413,1) (1.39329609,1) (1.39664804,1) (1.4,1)};\end{axis}\end{tikzpicture}\end{center}

这是题给平均场接触模型的失稳，不是实际电子关联体系的精确相图。
\end{solution}
\begin{exercise}[1999题2(a)--(g)]\label{f:333-1999-2}
spin1理想Bose气$\epsilon_{k,s}=\hbar^2k^2/(2m)-\mu_0sB$，$s=-1,0,+1$，忽略轨道磁性。(a)占据；(b)各$N_s$；(c)$M,\chi(T,\mu)$；(d)零场稀薄$z(n)$至$n^2$与磁化率修正；(e)$T_c(n,0)$和上侧$\chi$；(f)凝聚以下小$B$的$\mu$及凝聚态；(g)$B\downarrow0$自发磁化。
\end{exercise}
\begin{solution}
$z_s=z\e^{s\beta\mu_0B}$，$\avg{n_{k,s}}=[z_s^{-1}\e^{\beta\epsilon_k}-1]^{-1}$，$N_s=V\lambda^{-3}\operatorname{Li}_{3/2}(z_s)$（凝聚时另加最低态）。$M=\mu_0(N_+-N_-)$，$\chi_{B=0}=2\beta\mu_0^2V\lambda^{-3}\operatorname{Li}_{1/2}(z)$。

$x=n\lambda^3/3$，数式$x=z+z^2/2^{3/2}+\cdots$给$z=x-x^2/2^{3/2}$，所以$\chi=2N\mu_0^2[1+x/2^{3/2}+\cdots]/(3\kb T)$。$T_c=(2\pi\hbar^2/m\kb)[n/(3\zeta(3/2))]^{2/3}$，$z\to1$时$\operatorname{Li}_{1/2}$发散；数式$\operatorname{Li}_{3/2}(z)=\zeta(3/2)-2\sqrt\pi\sqrt{-\log z}+\cdots$给固定密度$\chi\propto(T-T_c)^{-1}$。

$B>0$最低态为$k=0,s=+1$，凝聚$\mu=-\mu_0B$。先热力学极限再$B\downarrow0$，三个激发组相同而凝聚全在$+$，$M=\mu_0N[1-(T/T_c)^{3/2}]$。严格有限体系先置$B=0$的对称混合给$M=0$；极限次序和固定粒子数条件必须注明。
\end{solution}
1999题3(a)--(d)全部为习题\ref{r:333-3-10}的经典/量子virial、星系估算，条件及每个小问相同。
\originalsection{2000年：氦3熔化与Bose占据}
\begin{exercise}[2000题1(a)--(f)]\label{f:333-2000-1}
He3熔化线低于0.3 K斜率负，0.1 K约$-30$ atm/K。(a)固相每核spin1/2、忽略核自旋作用的熵；(b)液相理想Fermi气每原子体积46 Å$^3$，估$T_F$；(c)低温热容；(d)熵与两相比较；(e)由$\mu$相等证明CC式；(f)$v_\ell-v_s=3$ Å$^3$每原子时低温斜率。
\end{exercise}
\begin{solution}
无耦合核spin各2态，$s_s=\kb\log2$。液体$n=1/(46\times10^{-30}\mathrm{m^3})$，裸He3质量约$5.0\times10^{-27}$ kg给$T_F=\hbar^2(3\pi^2n)^{2/3}/(2m\kb)\simeq6.1$ K。$C_V/N=(\pi^2/2)\kb T/T_F$、$s_\ell=(\pi^2/2)\kb T/T_F$，足够低温液体熵小于核自旋未冻结的固体。

$\dd\mu=-s\dd T+v\dd P$、两相$\dd\mu$相等，$\dd P/\dd T=(s_\ell-s_s)/(v_\ell-v_s)$。在所述低温窗口忽略液体熵，斜率$-\kb\log2/(3\times10^{-30}\mathrm{m^3})\simeq-3.19\times10^6$ Pa/K，约$-31.5$ atm/K，与题给尺度相合。核交换尺度以下固相spin也不能保持$\kb\log2$，故此“低温”是低于$T_F$且高于极小spin作用尺度，不可用于真实严格$T=0$。
\end{solution}
\begin{exercise}[2000题2(a)--(f)]\label{f:333-2000-2}
理想Bose单粒子态$q$、能$E_q$、$z_q=\e^{\beta(\mu-E_q)}<1$。(a)所有$n_q$的联合概率；(b)一模特征函数；(c),(d)均值、方差及方差对均值；(e)以均占据重写联合概率；(f)熵与零温。
\end{exercise}
\begin{solution}
每模几何分布，$P(\{n_q\})=\prod_q(1-z_q)z_q^{n_q}$、$\varphi_q(k)=(1-z_q)/(1-z_q\e^{\ii k})$，导数得$\bar n_q=z_q/(1-z_q)$、$\operatorname{Var}n_q=z_q/(1-z_q)^2=\bar n_q(1+\bar n_q)$。改写$P=\prod_q\bar n_q^{n_q}/(1+\bar n_q)^{1+n_q}$，熵
\begin{equation}
S=\kb\sum_q[(1+\bar n_q)\log(1+\bar n_q)-\bar n_q\log\bar n_q].
\end{equation}
无宏观激发时体熵趋0；凝聚基态的巨正则几何数涨落可仍巨大，单模熵约$\kb\log N_0$而每粒子熵0。固定总$N$的纯基态熵0。必须区分巨正则的基态粒子数不确定与有限温激发熵，不能由温度0无条件断言任何混合密度矩阵熵都0。
\end{solution}
2000题3(a)--(d)二维硬杆恰为既有热统习题33.5。临界角由等熵与驻值共同给，$\theta_c=1.19918$、$n_cl^2=0.265566$；驻值折点$1.72067,0.239624$为亚稳界，不误作全局转变。
\originalsection{2003、2006年：氦4；2003、2005年：Dirac激发}
\begin{exercise}[2003题1及2006题1，均(a)--(e)]\label{f:333-he4}
He4低温压致固化，熔化压力在0.8 K附近比0 K低约20 Pa。(a)CC式；(b)液相一支声速$c$的声子热容；(c)固相纵$c_L$及两横$c_T$的热容；(d)$c\simeq c_L\simeq c_T$、$v_\ell\simeq v_s=v$时熵差与哪个熵高；(e)恒小$\delta v=v_\ell-v_s$下熔化曲线与较高密度相。
\end{exercise}
\begin{solution}
CC式已在前题从$\mu$证明。每支三维声子态密度$V\omega^2/(2\pi^2c_s^3)$，$U=\int\hbar\omega\rho/(\e^{\beta\hbar\omega}-1)\dd\omega$给$C/N=2\pi^2v\kb^4T^3/(15\hbar^3c_s^3)$。固相将$c_s^{-3}$换成$c_L^{-3}+2c_T^{-3}$；熵积分$C/T$把系数除3再乘$T^3$。

若各声速和体积近同，记单支热容系数$a=2\pi^2v\kb^4/(15\hbar^3c^3)$，$s_\ell-s_s=-2aT^3/3$，固相比液体声子支更多、熵较高。因此$P(T)=P(0)-aT^4/(6\delta v)$。题述初始下降要求$\delta v>0$，固体体积小、密度高。该纯声子模型只解释低温下降趋势，不预测曲线在0.8 K反转的全部机制，也不在未给$c,v,\delta v$时伪算20 Pa幅度。
\end{solution}
2003题2(a)--(f)表面活性剂与既有习题33.1完全相同，原条件、完整解、塌缩/相分离近似及定张力热容均在那里。
\begin{exercise}[2003题3(a)--(f)；2005题1(a)--(e)]\label{f:333-dirac}
2003三维spin1/2 Dirac谱$E_\pm(k)=\pm\sqrt{m^2c^4+\hbar^2c^2k^2}$，零温负能满、正能空。(a)占据$\mu+\delta$与空穴$\mu-\delta$概率；(b)$\mu(T)$；(c)粒子空穴能积分；(d),(e)无质量能与热容；(f)有质量低温。2005石墨烯二维谱$E_\pm=\pm\hbar v|k|$、题给自旋$g=2,v=10^6$ m/s：(a)--(d)重复二维能和热容；(e)声速$2\times10^4$ m/s的声子与电子比较。按该题印出的$g=2$计，不再擅加另一个谷因子。
\end{exercise}
\begin{solution}
Fermi函数$f(\mu+\delta)=1-f(\mu-\delta)$，对称谱固定半填充$\mu=0$。粒子空穴因子2、自旋$g=2$，$\Delta U=2gV\int\dd^3k/(2\pi)^3\,E_+/(\e^{\beta E_+}+1)$。无质量径向积分$\int_0^\infty x^3/(\e^x+1)\dd x=7\pi^4/120$，
\begin{equation}
\Delta U/V=\frac{7g\pi^2}{120}\frac{(\kb T)^4}{\hbar^3c^3},\quad C_V/V=\frac{7g\pi^2}{30}\frac{\kb^4T^3}{\hbar^3c^3}.
\end{equation}
有隙$m c^2=\Delta$的低温粒子数$\propto T^{3/2}\e^{-\Delta/\kb T}$，能领先$\Delta$乘粒子数，热容$\propto T^{-1/2}\e^{-\Delta/\kb T}$（其余常数按近底有效质量），所以指数冻结而非无质量$T^3$。

二维$\Delta U/A=3g\zeta(3)(\kb T)^3/(2\pi\hbar^2v^2)$，$C/A=9g\zeta(3)\kb^3T^2/(2\pi\hbar^2v^2)$，与习题\ref{p:333-6-5}单层公式在$g=2,b=\hbar v$一致。线性声子也$T^2$但系数$\propto c_s^{-2}$，$v/c_s\sim50$使声子通常大许多；自由薄层还可有二次弯曲支。并非只凭电子无质量就判定电子控制总热容。
\end{solution}
\originalsection{2004年：零能壳与相对论Bose气}
2004题1(a)--(d)二次零能壳谱是习题\ref{p:333-6-4}，新增(a)化学势由上一题$f(E)+f(-E)=1$给$\mu=0$，其余积分、闭式与$\sqrt T$热容逐项在原题处。2004题2(a)--(d)三阶气体临界行为完整对应习题\ref{p:334-1-1}(a)--(d)。
\begin{exercise}[2004题3(a)--(e)]\label{f:333-2004-3}
$d$维守粒子数、无自旋Bose气$\epsilon=c|p|$。(a)巨势、数密度用$\operatorname{Li}_s(z)$；(b)$P,U,U/(PV)$与经典；(c)$T_c(n)$及维数；(d)凝聚以下热容温变；(e)$T_c$热容与经典极限，区分需要的一侧极限。
\end{exercise}
\begin{solution}
态密度$\rho(E)/V=S_{d-1}E^{d-1}/[(2\pi\hbar c)^d]$，定义$C_d=S_{d-1}\Gamma(d)/(2\pi\hbar c)^d$。能量积分或巨势分部积分给
\begin{equation}
\mathcal G/V=-C_d(\kb T)^{d+1}\operatorname{Li}_{d+1}(z),\quad n=C_d(\kb T)^d\operatorname{Li}_d(z),\quad U=dPV.
\end{equation}
经典每粒子$U=d\kb T$，与量子$U/(PV)=d$同。$d>1$才有$\zeta(d)$有限，$\kb T_c=[n/(C_d\zeta(d))]^{1/d}$；这是守粒子数气，不是可任意创灭而$\mu=0$的平衡光子。

凝聚以下与正常以上固定粒子数的热容分别为
\begin{align}
\frac{C_<}{N\kb}&=d(d+1)\frac{\zeta(d+1)}{\zeta(d)}\left(\frac{T}{T_c}\right)^d,\\
\frac{C_>}{N\kb}&=d(d+1)\frac{\operatorname{Li}_{d+1}}{\operatorname{Li}_d}
-d^2\frac{\operatorname{Li}_d}{\operatorname{Li}_{d-1}}.
\end{align}
临界\textbf{左}值相对经典$dN\kb$为$(d+1)\zeta(d+1)/\zeta(d)$；$1<d\le2$左右同值，$d>2$右侧须减$d\zeta(d)/\zeta(d-1)$，有有限热容跃变。例如$d=3$两侧约$10.805N\kb$、$4.228N\kb$，经典$3N\kb$。原题若只写“在$T_c$的热容”，须明确取凝聚侧，不能把左右混为一个数。
\end{solution}
\originalsection{2005、2006年的其余编号}
2005题2(a)--(f)相对论Coulomb气完整对应习题\ref{r:333-3-9}，同一齐次Hamilton量的缩放不由密度本身破坏。2006题3(a)--(e)双层石墨烯完整对应习题\ref{p:333-6-5}(a)--(e)，包括粒子空穴积分与两层电子热容交叉尺度。
\begin{exercise}[2005题3(a)--(f)]\label{f:333-2005-3}
理想Fermi气巨系综、单态$k$能$E_k$。(a)联合占据概率；(b)用平均占据重写；(c)$l$种概率$p_j$熵及最大值；(d)Fermi占据熵、零温；(e)总$N$方差、零温；(f)由$\operatorname{Var}N=Nn\kb T\kappa_T$估金属零温压缩率，单位Å$^3$/eV。
\end{exercise}
\begin{solution}
每态Bernoulli，$p_k=f(E_k)$，$P(\{n_k\})=\prod_kp_k^{n_k}(1-p_k)^{1-n_k}$、$n_k=0,1$。任意离散熵$S=-\kb\sum p_j\log p_j\le\kb\log l$，由凹性或相对均匀分布的相对熵，均匀时取等号。因此$S_F=-\kb\sum_k[p_k\log p_k+(1-p_k)\log(1-p_k)]$。

零温非简并Fermi海每态0或1，熵与数方差0，$\operatorname{Var}N=\sum_kp_k(1-p_k)$。但低温方差$\simeq\kb T\rho(E_F)$，除$T$有有限响应：$\kappa_T=\rho(E_F)/(Vn^2)=3/(2nE_F)$（三维二次色散）。取示例$n=0.1$ Å$^{-3}$、裸电子自由模型$E_F\simeq7.86$ eV，得$1.91$ Å$^3$/eV。准确金属参数各不相同，原题未给$n$故估值必须列假设，不能把$T=0$直接代涨落式得到错误$0/0$。
\end{solution}
\begin{exercise}[2006题2(a)--(f)]\label{f:333-2006-2}
链$N$个长$a$无质量连杆在细管中只能平行/反平行，端距$L$，$x=L/(Na)$。(a)态数；(b)Stirling熵；压挤给原题近似能$E=-\sigma a^2Nx^2/2$，用于每个给定$L$的态。(c)自由能及端力$J$，功$J\dd L$；(d)高低温等温线与$T_c$；(e)稳定与不稳定范围；(f)无外力近$T_c$下长度。
\end{exercise}
\begin{solution}
$N_\pm=N(1\pm x)/2$，$\Omega=\binom N{N_+}$（奇偶允许），$S/(N\kb)=\log2-[(1+x)\log(1+x)+(1-x)\log(1-x)]/2$。$F=E-TS$，力$J=F_L=(\kb T/a)\operatorname{atanh}x-\sigma ax$。$\kb T_c=\sigma a^2$，小$x$力斜率过零；下图纵轴$J/(\sigma a)$。
\begin{center}\begin{tikzpicture}\begin{axis}[width=10cm,height=4.6cm,xlabel={$x$},ylabel={$J/(\sigma a)$},xmin=-1,xmax=1,ymin=-1.7,ymax=1.7,legend pos=north west]\addplot[structurecolor,domain=-.96:.96,samples=180]{1.3*.5*ln((1+x)/(1-x))-x};\addlegendentry{$T/T_c=1.3$}\addplot[gray,domain=-.96:.96,samples=180]{.5*ln((1+x)/(1-x))-x};\addlegendentry{$T/T_c=1$}\addplot[main,domain=-.96:.96,samples=180]{.6*.5*ln((1+x)/(1-x))-x};\addlegendentry{$T/T_c=.6$}\end{axis}\end{tikzpicture}\end{center}
局部稳定$F_{LL}=[\kb T/(1-x^2)-\sigma a^2]/(Na^2)>0$，低温$|x|<\sqrt{1-T/T_c}$内不稳定。全局平衡要比较两个极小/凸包，不把所有稳定支都当全局最优。无力$T/T_c=x/\operatorname{atanh}x$，展开$\operatorname{atanh}x=x+x^3/3+\cdots$，$|L|\simeq Na\sqrt{3(1-T/T_c)}$。高温$J=0$只有$x=0$，低温两种对称伸长方向形成题给平均场分岔。
\end{solution}
